\documentclass[11pt]{article}
\usepackage[margin=0.85in]{geometry}
\usepackage{lmodern}
\usepackage{microtype}
\usepackage{enumitem}
\usepackage{listings}
\usepackage[T1]{fontenc}

\lstset{
  basicstyle=\ttfamily\small,
  columns=fullflexible,
  keepspaces=true,
  showstringspaces=false,
  breaklines=true,
  frame=single,
  xleftmargin=0.4em,
  xrightmargin=0.4em,
  aboveskip=0.7em,
  belowskip=0.7em
}

\pagestyle{empty}
\setlength{\parskip}{0.45em}
\setlength{\parindent}{0pt}
\setlist[enumerate]{leftmargin=1.4em, itemsep=0.85em, topsep=0.4em}
\setlist[itemize]{leftmargin=1.3em, itemsep=0.25em, topsep=0.35em}

\begin{document}

\begin{center}
{\LARGE\bfseries Quiz 1 --- Sample 1}\\[0.35em]
DATA 412/612
\end{center}

\vspace{0.4em}

Write your answers in an R script. You may use R and RStudio.
You may \textbf{not} use the internet or your notes.
Do \textbf{not} use dplyr, \texttt{filter()}, or \texttt{select()} on this quiz.
Use Base R, including \texttt{subset()} and \texttt{function()}.

\begin{enumerate}
\item \textbf{\texttt{subset()}}

This data frame is already in your environment:

\begin{lstlisting}
offers <- data.frame(
  firm   = c("Helix", "Nimble", "Helix", "Orbit", "Nimble", "Orbit", "Helix"),
  role   = c("intern", "intern", "analyst", "intern", "analyst", "analyst", "intern"),
  salary = c(22, 19, NA, 21, 28, 27, 24),
  remote = c(TRUE, FALSE, TRUE, TRUE, FALSE, TRUE, TRUE)
)
\end{lstlisting}

\begin{itemize}
  \item Using \texttt{subset()} only (not \texttt{[} and not dplyr), keep the
        \textbf{intern} rows whose \texttt{salary} is \textbf{strictly greater than the
        mean salary of interns}. Compute that intern-only mean with missing values
        removed. Drop the \texttt{role} column from what you return.
  \item Using \texttt{subset()} on the \texttt{salary} \textbf{vector}
        (not on the whole data frame), return every salary that is at least 22.
\end{itemize}

\item \textbf{Writing a function}

Write a function named \texttt{pair\_mean\_call} in an R script. It takes one
argument, a data frame \texttt{df}.

\begin{itemize}
  \item If \texttt{df} has fewer than 2 columns, return \texttt{"thin"}.
  \item If \texttt{df} has more than 2 columns, return \texttt{"wide"}.
  \item If the first column or the second column is not numeric, return \texttt{FALSE}.
  \item If either of those two columns contains any \texttt{NA}, return \texttt{"missing"}.
  \item Using a pipe, compute the mean of the \textbf{finite} values in column 1
        and the mean of the \textbf{finite} values in column 2.
        (\texttt{Inf}, \texttt{-Inf}, and \texttt{NaN} are not finite.
        They are also not \texttt{NA}; do not treat them as missing for the
        previous bullet.)
  \item Return \texttt{1} if the first mean is greater than or equal to the
        second mean. Otherwise return the \textbf{name} of the second column.
\end{itemize}

Run your function on each of these inputs:

\begin{lstlisting}
data.frame(u = 1:2)
data.frame(u = 1:2, v = c("a", "b"))
data.frame(u = c(1, NA), v = c(2, 3))
data.frame(u = c(1, Inf, 3), v = c(10, 10, 10))
data.frame(low = c(2, 4), high = c(1, 1))
data.frame(low = c(1, 1), high = c(2, 4))
\end{lstlisting}
\end{enumerate}

\end{document}
